\subsubsection*{4.4. Griffiths 定理的证明}

设 $z$ 为 $X$ 上余维为 $m$ 的代数圈, 使得
\begin{equation}\tag{4.4.1}
z_{\bar\eta}=z|_{X_{\bar\eta}}
\quad\text{代数等价于零}.
\end{equation}
按定义, 存在
\begin{equation}\tag{4.4.2}
\text{$k(\bar\eta)$ 上固有、光滑且连通的曲线 $C_{\bar\eta}$},
\end{equation}
\begin{equation}\tag{4.4.3}
t_i(\bar\eta):\Spec k(\bar\eta)\longrightarrow C_{\bar\eta},
\qquad i=1,2,
\end{equation}
以及
\begin{equation}\tag{4.4.4}
\text{$C_{\bar\eta}\times X_{\bar\eta}$ 上余维为 $m$ 的代数圈 $P_{\bar\eta}$},
\end{equation}
使得
\begin{equation}\tag{4.4.5}
z_{\bar\eta}
=P_{\bar\eta}\bigl(t_1(\bar\eta)\bigr)
-P_{\bar\eta}\bigl(t_2(\bar\eta)\bigr)
\quad\text{于 }X_{\bar\eta}.
\end{equation}
印本在 (4.4.3) 中将目标写作 $C$ 而不是 $C_{\bar\eta}$, 并在上文以及 (4.4.4) 中将余维写作 $n$ 而不是 $m$; 所定义的曲线、定理~4.1 以及下文的次数移位均要求此处采用的读法.

将 $P_{\bar\eta}$ 视为从 $C_{\bar\eta}$ 到 $X_{\bar\eta}$ 的代数对应. 它定义一个仍记为 $P_{\bar\eta}$ 的上同调对应, 或等价地定义同态
\begin{equation}\tag{4.4.6}
P_{\bar\eta}:H^*(C_{\bar\eta},\mathbb Q_\ell(i))
\longrightarrow
H^{*+2(m-1)}(X_{\bar\eta},\mathbb Q_\ell(i+m-1)).
\end{equation}
印本在 (4.4.6) 的目标中漏去了最后一个右括号.

同样地, 考虑 (4.1.3) 中假设 (B) 所给出的 (3.3.2) 中从 $X_{\bar\eta}$ 到 $X\otimes_k k(\bar\eta)$ 的代数对应 $V_{\bar\eta}$. 它定义
\begin{equation}\tag{4.4.7}
V_{\bar\eta}:H^*(X_{\bar\eta},\mathbb Q_\ell(i+m-1))
\longrightarrow
H^*(X\otimes_k k(\bar\eta),\mathbb Q_\ell(i+m-1)).
\end{equation}
因此, 从 $C_{\bar\eta}$ 到 $X\otimes_k k(\bar\eta)$ 的复合代数对应 $Q_{\bar\eta}=V_{\bar\eta}\circ P_{\bar\eta}$ 定义
\begin{equation}\tag{4.4.8}
Q_{\bar\eta}=V_{\bar\eta}\circ P_{\bar\eta}:
H^*(C_{\bar\eta},\mathbb Q_\ell(i))
\longrightarrow
H^{*+2(m-1)}(X\otimes_k k(\bar\eta),\mathbb Q_\ell(i+m-1)).
\end{equation}

考虑包含映射
\begin{equation*}
m:X_{\bar\eta}\longrightarrow X\otimes_k k(\bar\eta).
\end{equation*}
它给出从 $X\otimes_k k(\bar\eta)$ 到 $X_{\bar\eta}$ 的代数对应 $m^*$, 从而给出从 $C_{\bar\eta}$ 到 $X_{\bar\eta}$ 的复合对应 $m^*Q_{\bar\eta}$, 定义
\begin{equation}\tag{4.4.9}
m^*Q_{\bar\eta}=m^*V_{\bar\eta}P_{\bar\eta}:
H^*(C_{\bar\eta},\mathbb Q_\ell(i))
\longrightarrow
H^{*+2(m-1)}(X_{\bar\eta},\mathbb Q_\ell(i+m-1)).
\end{equation}
在上述句子中, 印本在 $m^*Q_{\bar\eta}$ 之后重复写了一个多余的 $\bar\eta$.

现在考虑代数对应
\begin{equation}\tag{4.4.10}
\begin{split}
R_{\bar\eta}
&=P_{\bar\eta}-m^*Q_{\bar\eta}
=(\id-m^*V_{\bar\eta})P_{\bar\eta}:\\
&\hspace{2em}H^*(C_{\bar\eta},\mathbb Q_\ell(i))
\longrightarrow
H^{*+2(m-1)}(X_{\bar\eta},\mathbb Q_\ell(i+m-1)).
\end{split}
\end{equation}
在 $n-1=2m-1$ 次上同调上, $V_{\bar\eta}$ 是 $m^*$ 的左逆, 因而 $\id-m^*V_{\bar\eta}$ 是 $H^{n-1}(X_{\bar\eta},-)$ 上的投影算子, 其像为
\begin{equation*}
\operatorname{Im}\bigl(m^*:H^{n-1}(X\otimes_k k(\bar\eta),-)
\longrightarrow H^{n-1}(X_{\bar\eta},-)\bigr)
\end{equation*}
的一个补空间.\EditorialFootnote{译校: 法文印本第 358 页将本句的上同调次数统写为 $*$; 条件 (B), 即 (3.3.3)--(3.3.5), 只保证 $n-1$ 次上的左逆关系. 这里及下文 (4.4.11) 只需 $n-1=2m-1$ 次, 故明确限定于此次数.} 由假设 (4.1.2),
\begin{equation}\tag{4.4.11}
R_{\bar\eta}:H^1(C_{\bar\eta},\mathbb Q_\ell(i))
\longrightarrow
H^{2m-1}(X_{\bar\eta},\mathbb Q_\ell(i+m-1))
\end{equation}
的像位于 $\operatorname{Im}m^*$ 中. 所以,
\begin{equation}\tag{4.4.12}
\text{同态 (4.4.11) 为零.}
\end{equation}

选择有限扩张
\begin{equation*}
L/k(\eta)
\end{equation*}
使得对象
\begin{equation*}
C_{\bar\eta},\quad t_1(\bar\eta),\quad t_2(\bar\eta),\quad
Q_{\bar\eta},\quad R_{\bar\eta}
\end{equation*}
均已在其上有定义, 并令 $U$ 为 $\mathbb P^1$ 在 $L$ 中的正规化的非空开子集. 缩小 $U$ 后, 存在纤维几何连通、几何泛纤维为 $C_{\bar\eta}$ 的固有光滑相对曲线
\begin{equation*}
C\longrightarrow U
\end{equation*}
及两个截面
\begin{equation*}
t_i:U\longrightarrow C,\qquad i=1,2,
\end{equation*}
它们在 $\bar\eta$ 处的值为 $t_i(\bar\eta)$.

设
\begin{equation*}
\rho:\widetilde X\longrightarrow\mathbb P^1
\end{equation*}
为 XVIII~3.1.4 中以束的各成员为纤维的投影, 置
\begin{equation*}
\widetilde X_U=\widetilde X\times_{\mathbb P^1}U,
\end{equation*}
并设
\begin{equation*}
M:\widetilde X_U\hookrightarrow X\times U
\end{equation*}
为到 $U$ 上纤维为 $X$ 的常族的典范浸入. 置
\begin{equation*}
z_U=M^*\pr_1^*(z),
\end{equation*}
即由 $z$ 诱导的 $\widetilde X_U$ 上的代数圈.

再次缩小 $U$ 后, 存在 $C\times_U\widetilde X_U$ 上的代数圈 $R$, 它在 $C_{\bar\eta}\times X_{\bar\eta}$ 上的值为 $R_{\bar\eta}$, 以及 $C\times X\times U$ 上的代数圈 $Q$, 它在 $C_{\bar\eta}\times(X\otimes_k k(\bar\eta))$ 上的值为 $Q_{\bar\eta}$, 使得
\begin{equation}\tag{4.4.13}
z_U=R(t_1)-R(t_2)
+(\id\times M)^*\bigl[Q(t_1)-Q(t_2)\bigr].
\end{equation}

由 Griffiths 同态``模去固定部分'' (XVIII~5.8.7) 后的单射性,
\begin{equation*}
\operatorname{Prim}^{2m}(X,\mathbb Q_\ell(m))
\xrightarrow{\ \underline{\operatorname{grif}}\ }
\frac{H^1(U,R^{2m-1}\rho_{U*}\mathbb Q_\ell(m))}
{M^*H^1(U,H^{2m-1}(X,\mathbb Q_\ell(m)))},
\end{equation*}
为证明定理, 只须证明 $z_U$ 经
\begin{equation*}
u_0:\operatorname{Prim}^{2m}(\widetilde X_U/U,\mathbb Q_\ell(m))
\longrightarrow H^1(U,R^{2m-1}\rho_{U*}\mathbb Q_\ell(m))
\end{equation*}
所得的像位于 $M^*H^1(U,H^{2m-1}(X,\mathbb Q_\ell(m)))$ 中.
印本将 Griffiths 映射定义域中的扭写作 $n$; 余维及所有相邻项均要求 $m$.

由 (4.4.13) 和 (1.4.2) 的交换性, 有交换图
\begin{center}
\begin{tikzcd}[column sep=4.2em,row sep=4.2em]
\operatorname{Prim}^{2}(C/U,\mathbb Q_\ell(1))
  \arrow[r,"Q"]
& \operatorname{Prim}^{2m}(X\times U/U,\mathbb Q_\ell(m))
  \arrow[r,"M^*"] \arrow[d,"u_{\pr_2}"]
& \operatorname{Prim}^{2m}(\widetilde X_U/U,\mathbb Q_\ell(m))
  \arrow[d,"u_0"] \\
{}
& H^1(U,H^{2m-1}(X,\mathbb Q_\ell(m)))
  \arrow[r,"M^*"]
& H^1(U,R^{2m-1}\rho_{U*}\mathbb Q_\ell(m)),
\end{tikzcd}
\end{center}
所以只须证明
\begin{equation*}
u_0\bigl(R(t_1)-R(t_2)\bigr)=0
\quad\text{于}\quad
H^1(U,R^{2m-1}\rho_{U*}\mathbb Q_\ell(m)).
\end{equation*}
由 1.7, 这从 (4.4.12) 得到. 印本在最后一群中写作 $2n-1$ 并使用扭 $n$, 还引用了不存在的 (4.3.12); 相容的读法是 $2m-1$、扭 $m$ 和 (4.4.12).
